\begin{abstract}
We study a recognition-dependent observer-selection framework motivated by Quantum Bogosort (QBS). Recognition changes an agent's policy, and the resulting policy may change both downstream trajectory utility and observer-indexed accessibility. For a base measure $\mu$, utility $U_\pi$, and nonnegative accessibility $S_\pi$, first-person value is defined by
\[
V_{\mathrm{FP}}(\pi)=\frac{\mathbb E_\mu[U_\pi S_\pi]}{\mathbb E_\mu[S_\pi]}.
\]
We prove covariance and tail-probability identities, a monotone-accessibility first-order stochastic dominance theorem, an exact recognition decomposition separating ordinary trajectory effects from first-person conditioning effects, and a policy--selection interaction decomposition. Classical simulations test the theorem boundaries, show that predictive alignment can arise endogenously in a minimal learned agent, and distinguish marginal first-person uplift from cross-branch decision correlation. The mathematical results are separated from a distinct Everett-QBS bridge assumption identifying accessibility weighting with first-person self-location. The present work therefore establishes a formal decision-and-conditioning framework, not a derivation of observer weighting from Everettian quantum mechanics.
\end{abstract}
